\Needspace{12\baselineskip}
\section{总结}
\subsection{可扩展之处}
当前主工程已统一为\path{camera processing unit/}目录结构，主板采用WBCAtt+与BCCD两来源重训模型。SRAM下载已核实，用户反馈实拍改善并自行完成Flash烧录；手动烧录文件身份与冷启动尚未独立核验。此前单一来源新版因空场景误检撤回，历史证据保留。系统实现全帧量化联合检测分类、15 Hz抽帧、单框身份保持及英文全称。后续先采集主板相机带标签细胞数据，检查淋巴细胞弱项、照明与色彩偏差及小目标细节；再考虑输入分辨率、真正多尺度检测或多目标计数。剩余14块ERAM与5个DSP，扩大网络需综合吞吐、资源和时序。新增实时功能继续在板端实现。

\subsection{心得体会}
本项目的核心是把图像、时序和控制放在同一套可核对的数据约定下。RAW10 的组宽、有效行数、Bayer 相位、窗口中心和网络偏移都会影响最终画面的空间位置。只有在帧首、有效像素和坐标之间建立一致关系，算法模块才能在不改变主链路的条件下连接和复用。

自动图像调节同样需要明确反馈来源和执行时刻。AWB 应估计数字校正前的色偏，AE 应测量校正后亮度；曝光寄存器写入完成与图像反映新曝光之间存在延迟。控制器通过窗口平均、死区、限幅、积分管理和生效等待，把这些实际约束纳入硬件决策，参数提交则以帧边界保证整帧一致。

网络显示不是把 UDP 字节直接写进显存即可完成。包到达顺序、缺口填充、FIFO 门槛、DMA 突发和完整帧发布必须协同。数据从 DDR 返回时，采样基准应与 DQS 对应，跨域后的有效信号应描述真实收到的数据；进入 HDMI 前还需把 RGB 和同步信号延迟一致。这些约束共同决定图像是否在正确位置连续显示。

验证应覆盖接口两端和异常路径。像素参考模型用于核对数值，网络检查器用于核对线上格式，生产存储控制器和芯片模型用于核对读写关系，显示端逐像素比对用于核对最终光栅。功能仿真、静态时序和上板观察分别回答不同问题，将它们结合，才能形成可持续扩展的实时视觉系统。

检测分类接入表明，数据定义、预处理、量化和展示策略需要一起验证。WBCAtt+掩膜提供整细胞框和五类，BCCD只监督定位；去重与冲突剔除应在性能查看前固定。本轮权重由VAL选择，部署阈值另用开发截图校准，这些截图不能再作为独立验收。整数逐层对拍回答计算是否一致，完整视频仿真回答调度和结果保持是否正确，物理时序回答既有约束下能否运行。用户实拍改善反馈提供了定性证据，后续仍需独立带标签相机集、真实空场景统计与持续运行检查，才能量化公开数据结果在本板拍摄域中的表现。
